\documentclass[12pt,a4paper]{article}
\usepackage[utf8]{vietnam}
\usepackage{amsmath,amssymb}
\usepackage{tabularx}
\usepackage{tikz}
\usepackage{xcolor}
\usepackage[top=0.6cm, bottom=0.6cm, left=0.6cm, right=0.6cm]{geometry}
\usepackage[dethi]{ex_test}

\SAOPTN{widthSA=3.2,heightSA=0.62}

\setlength{\parskip}{0pt}

% Dòng kẻ chấm mờ đều đặn chuẩn viết tay
\newcommand{\dongke}[1]{%
  \begingroup%
  \setlength{\parskip}{0pt}%
  \baselineskip=18.5pt%
  \vspace{1.0mm}%
  \foreach \i in {1,...,#1}{%
    \noindent{\color{gray!60}\dotfill}\par%
  }%
  \vspace{0.6mm}%
  \endgroup%
}

% Định nghĩa ô bài tập tự luận trong bảng
\newcommand{\oBaiTap}[3][6]{%
  \vspace{1.0mm}%
  \noindent\textbf{#2} #3\par%
  \dongke{#1}%
}

\begin{document}
\pagestyle{empty}
% =============================================================================
% ========================= TOÁN 9 - TUẦN 06 - BUỔI 01: MÃ ĐỀ 101 =============
% ========================= MẶT TRƯỚC (TRANG 1/2) ==============================
% =============================================================================

\noindent
\begin{tabularx}{\linewidth}{|X|c|}
\hline
\textbf{PHIẾU KIỂM TRA ĐỊNH KỲ TOÁN 9 - TUẦN 06 - BUỔI 01} & \textbf{Mã đề: 101 - Trang 1/2} \\
\hline
\end{tabularx}

\vspace{1.0mm}

\noindent
\begin{tabularx}{\linewidth}{|p{3.2cm}|X|}
\hline
\centering\textbf{ĐIỂM SỐ} & \centering\textbf{LỜI NHẬN XÉT CỦA THẦY CÔ} \tabularnewline
\hline
\rule{0pt}{1.8cm} & \tabularnewline
\hline
\end{tabularx}

\vspace{1.2mm}
\noindent\textbf{Câu 1 (0,25 điểm).} Cho $\Delta ABC$ vuông tại $A$ có $AB = 3, BC = 5$. Khi đó $\sin B =$ \shortans{\dfrac{4}{5}}\par\vspace{0.8mm}
\noindent\textbf{Câu 2 (0,25 điểm).} Cho $\Delta ABC$ vuông tại $A$ có $AB = 3, BC = 5$. Khi đó $\cos B =$ \shortans{\dfrac{3}{5}}\par\vspace{0.8mm}
\noindent\textbf{Câu 3 (0,25 điểm).} Cho $\Delta ABC$ vuông tại $A$ có $AB = 4, AC = 3$. Khi đó $\tan B =$ \shortans{\dfrac{3}{4}}\par\vspace{0.8mm}
\noindent\textbf{Câu 4 (0,25 điểm).} Cho $\Delta ABC$ vuông tại $A$, đường cao $AH$. Biết $BH=4, CH=9$. Độ dài $AH$ là: \shortans{6}\par\vspace{0.8mm}
\noindent\textbf{Câu 5 (0,25 điểm).} Một cột cờ cao $6\text{ m}$ có bóng trên mặt đất dài $6\text{ m}$. Góc tạo bởi tia nắng và mặt đất là: \shortans{45^\circ}\par\vspace{0.8mm}
\noindent\textbf{Câu 6 (0,25 điểm).} Cặp $(1; y)$ là nghiệm của phương trình $2x + y = 5$. Khi đó $y =$ \shortans{3}\par\vspace{0.8mm}
\noindent\textbf{Câu 7 (0,25 điểm).} Tìm $m$ để cặp $(1; 2)$ là nghiệm phương trình $mx + y = 5$: \shortans{3}\par\vspace{0.8mm}
\noindent\textbf{Câu 8 (0,25 điểm).} Tìm số thực $x \ge 0$, biết $\sqrt{x} = 5$: \shortans{25}\par\vspace{0.8mm}
\noindent\textbf{Câu 9 (0,25 điểm).} Nghiệm nguyên dương của phương trình $(x - 3)(x + 2) = 0$ là: \shortans{3}\par\vspace{0.8mm}
\noindent\textbf{Câu 10 (0,25 điểm).} Nghiệm của phương trình $\sqrt{2x + 1} = 3$ là $x =$ \shortans{4}\par\vspace{0.8mm}
\vspace{0.8mm}
\noindent\textbf{Câu 11 (2,5 điểm). Giải các hệ phương trình sau:}
\vspace{0.6mm}

\noindent
\begin{tabularx}{\linewidth}{|X|X|}
\hline
\oBaiTap[21]{a) (1,25 điểm)}{$\begin{cases} 2x + y = 7 \\ x - y = 2 \end{cases}$} & 
\oBaiTap[21]{b) (1,25 điểm)}{$\begin{cases} 2(x - 1) + y = 4 \\ 3x - 2y = 2 \end{cases}$} \\
\hline
\end{tabularx}

\newpage
% =============================================================================
% ========================= TOÁN 9 - TUẦN 06 - BUỔI 01: MÃ ĐỀ 101 =============
% ========================= MẶT SAU (TRANG 2/2) ================================
% =============================================================================

\noindent
\begin{tabularx}{\linewidth}{|X|c|}
\hline
\textbf{TOÁN 9 - TUẦN 06 - BUỔI 01} & \textbf{Mã đề: 101 - Trang 2/2} \\
\hline
\end{tabularx}

\vspace{1.2mm}

\noindent\textbf{Câu 12 (2,5 điểm). Biến đổi và rút gọn biểu thức:}
\vspace{0.6mm}

\noindent
\begin{tabularx}{\linewidth}{|X|X|}
\hline
\oBaiTap[13]{a) (1,25 điểm)}{Rút gọn với $x > 0, x \ne 9$:\par $A = \dfrac{\sqrt{x}}{\sqrt{x} - 3} - \dfrac{9}{x - 3\sqrt{x}}$} & 
\oBaiTap[13]{b) (1,25 điểm)}{Rút gọn với $x > 0, x \ne 4$:\par $B = \left(\dfrac{1}{\sqrt{x} - 2} + \dfrac{1}{\sqrt{x} + 2}\right) \cdot \dfrac{x - 4}{2\sqrt{x}}$} \\
\hline
\end{tabularx}

\vspace{2.5mm}

\noindent
\begin{tabularx}{\linewidth}{|X|}
\hline
\vspace{0.5mm}
\textbf{Câu 13 (2,5 điểm).} Hai bạn An và Bình đến một nhà sách mua đồ dùng học tập. Bạn An mua $3$ quyển vở và $2$ chiếc bút bi cùng loại hết tổng cộng $34\,000$ đồng. Bạn Bình mua $2$ quyển vở và $4$ chiếc bút bi cùng loại hết tổng cộng $36\,000$ đồng. Tính giá tiền của mỗi quyển vở và mỗi chiếc bút bi.\par
\vspace{0.4mm}
\noindent\textit{Bài giải:}\dongke{22}
\\ \hline
\end{tabularx}

\newpage
% =============================================================================
% ========================= TOÁN 9 - TUẦN 06 - BUỔI 01: MÃ ĐỀ 102 =============
% ========================= MẶT TRƯỚC (TRANG 1/2) ==============================
% =============================================================================

\noindent
\begin{tabularx}{\linewidth}{|X|c|}
\hline
\textbf{PHIẾU KIỂM TRA ĐỊNH KỲ TOÁN 9 - TUẦN 06 - BUỔI 01} & \textbf{Mã đề: 102 - Trang 1/2} \\
\hline
\end{tabularx}

\vspace{1.0mm}

\noindent
\begin{tabularx}{\linewidth}{|p{3.2cm}|X|}
\hline
\centering\textbf{ĐIỂM SỐ} & \centering\textbf{LỜI NHẬN XÉT CỦA THẦY CÔ} \tabularnewline
\hline
\rule{0pt}{1.8cm} & \tabularnewline
\hline
\end{tabularx}

\vspace{1.2mm}
\noindent\textbf{Câu 1 (0,25 điểm).} Cho $\Delta ABC$ vuông tại $A$ có $AB = 6, BC = 10$. Khi đó $\sin B =$ \shortans{\dfrac{4}{5}}\par\vspace{0.8mm}
\noindent\textbf{Câu 2 (0,25 điểm).} Cho $\Delta ABC$ vuông tại $A$ có $AB = 6, BC = 10$. Khi đó $\cos B =$ \shortans{\dfrac{3}{5}}\par\vspace{0.8mm}
\noindent\textbf{Câu 3 (0,25 điểm).} Cho $\Delta ABC$ vuông tại $A$ có $AB = 8, AC = 6$. Khi đó $\tan B =$ \shortans{\dfrac{3}{4}}\par\vspace{0.8mm}
\noindent\textbf{Câu 4 (0,25 điểm).} Cho $\Delta ABC$ vuông tại $A$, đường cao $AH$. Biết $BH=1, CH=4$. Độ dài $AH$ là: \shortans{2}\par\vspace{0.8mm}
\noindent\textbf{Câu 5 (0,25 điểm).} Một cột đèn cao $5\text{ m}$ có bóng trên mặt đất dài $5\text{ m}$. Góc tạo bởi tia nắng và mặt đất là: \shortans{45^\circ}\par\vspace{0.8mm}
\noindent\textbf{Câu 6 (0,25 điểm).} Cặp $(2; y)$ là nghiệm của phương trình $3x + y = 8$. Khi đó $y =$ \shortans{2}\par\vspace{0.8mm}
\noindent\textbf{Câu 7 (0,25 điểm).} Tìm $m$ để cặp $(2; 1)$ là nghiệm phương trình $mx + 2y = 8$: \shortans{3}\par\vspace{0.8mm}
\noindent\textbf{Câu 8 (0,25 điểm).} Tìm số thực $x \ge 0$, biết $\sqrt{x} = 6$: \shortans{36}\par\vspace{0.8mm}
\noindent\textbf{Câu 9 (0,25 điểm).} Nghiệm nguyên dương của phương trình $(x - 5)(x + 1) = 0$ là: \shortans{5}\par\vspace{0.8mm}
\noindent\textbf{Câu 10 (0,25 điểm).} Nghiệm của phương trình $\sqrt{3x + 1} = 4$ là $x =$ \shortans{5}\par\vspace{0.8mm}
\vspace{0.8mm}
\noindent\textbf{Câu 11 (2,5 điểm). Giải các hệ phương trình sau:}
\vspace{0.6mm}

\noindent
\begin{tabularx}{\linewidth}{|X|X|}
\hline
\oBaiTap[21]{a) (1,25 điểm)}{$\begin{cases} 3x + y = 8 \\ x - y = 0 \end{cases}$} & 
\oBaiTap[21]{b) (1,25 điểm)}{$\begin{cases} 3(x + y) - 2x = 7 \\ 2x - y = 0 \end{cases}$} \\
\hline
\end{tabularx}

\newpage
% =============================================================================
% ========================= TOÁN 9 - TUẦN 06 - BUỔI 01: MÃ ĐỀ 102 =============
% ========================= MẶT SAU (TRANG 2/2) ================================
% =============================================================================

\noindent
\begin{tabularx}{\linewidth}{|X|c|}
\hline
\textbf{TOÁN 9 - TUẦN 06 - BUỔI 01} & \textbf{Mã đề: 102 - Trang 2/2} \\
\hline
\end{tabularx}

\vspace{1.2mm}

\noindent\textbf{Câu 12 (2,5 điểm). Biến đổi và rút gọn biểu thức:}
\vspace{0.6mm}

\noindent
\begin{tabularx}{\linewidth}{|X|X|}
\hline
\oBaiTap[13]{a) (1,25 điểm)}{Rút gọn với $x > 0, x \ne 4$:\par $A = \dfrac{\sqrt{x}}{\sqrt{x} - 2} - \dfrac{4}{x - 2\sqrt{x}}$} & 
\oBaiTap[13]{b) (1,25 điểm)}{Rút gọn với $x > 0, x \ne 9$:\par $B = \left(\dfrac{1}{\sqrt{x} - 3} + \dfrac{1}{\sqrt{x} + 3}\right) \cdot \dfrac{x - 9}{2\sqrt{x}}$} \\
\hline
\end{tabularx}

\vspace{2.5mm}

\noindent
\begin{tabularx}{\linewidth}{|X|}
\hline
\vspace{0.5mm}
\textbf{Câu 13 (2,5 điểm).} Hai bạn An và Bình đến một nhà sách mua đồ dùng học tập. Bạn An mua $4$ quyển vở và $2$ chiếc bút bi cùng loại hết tổng cộng $42\,000$ đồng. Bạn Bình mua $3$ quyển vở và $5$ chiếc bút bi cùng loại hết tổng cộng $49\,000$ đồng. Tính giá tiền của mỗi quyển vở và mỗi chiếc bút bi.\par
\vspace{0.4mm}
\noindent\textit{Bài giải:}\dongke{22}
\\ \hline
\end{tabularx}

\newpage
% =============================================================================
% ========================= TOÁN 9 - TUẦN 06 - BUỔI 01: MÃ ĐỀ 103 =============
% ========================= MẶT TRƯỚC (TRANG 1/2) ==============================
% =============================================================================

\noindent
\begin{tabularx}{\linewidth}{|X|c|}
\hline
\textbf{PHIẾU KIỂM TRA ĐỊNH KỲ TOÁN 9 - TUẦN 06 - BUỔI 01} & \textbf{Mã đề: 103 - Trang 1/2} \\
\hline
\end{tabularx}

\vspace{1.0mm}

\noindent
\begin{tabularx}{\linewidth}{|p{3.2cm}|X|}
\hline
\centering\textbf{ĐIỂM SỐ} & \centering\textbf{LỜI NHẬN XÉT CỦA THẦY CÔ} \tabularnewline
\hline
\rule{0pt}{1.8cm} & \tabularnewline
\hline
\end{tabularx}

\vspace{1.2mm}
\noindent\textbf{Câu 1 (0,25 điểm).} Cho $\Delta ABC$ vuông tại $A$ có $AB = 5, BC = 13$. Khi đó $\sin B =$ \shortans{\dfrac{12}{13}}\par\vspace{0.8mm}
\noindent\textbf{Câu 2 (0,25 điểm).} Cho $\Delta ABC$ vuông tại $A$ có $AB = 5, BC = 13$. Khi đó $\cos B =$ \shortans{\dfrac{5}{13}}\par\vspace{0.8mm}
\noindent\textbf{Câu 3 (0,25 điểm).} Cho $\Delta ABC$ vuông tại $A$ có $AB = 12, AC = 5$. Khi đó $\tan B =$ \shortans{\dfrac{5}{12}}\par\vspace{0.8mm}
\noindent\textbf{Câu 4 (0,25 điểm).} Cho $\Delta ABC$ vuông tại $A$, đường cao $AH$. Biết $BH=2, CH=8$. Độ dài $AH$ là: \shortans{4}\par\vspace{0.8mm}
\noindent\textbf{Câu 5 (0,25 điểm).} Một cây cau cao $8\text{ m}$ có bóng trên mặt đất dài $8\text{ m}$. Góc tạo bởi tia nắng và mặt đất là: \shortans{45^\circ}\par\vspace{0.8mm}
\noindent\textbf{Câu 6 (0,25 điểm).} Cặp $(1; y)$ là nghiệm của phương trình $4x + y = 7$. Khi đó $y =$ \shortans{3}\par\vspace{0.8mm}
\noindent\textbf{Câu 7 (0,25 điểm).} Tìm $m$ để cặp $(1; 3)$ là nghiệm phương trình $mx + y = 7$: \shortans{4}\par\vspace{0.8mm}
\noindent\textbf{Câu 8 (0,25 điểm).} Tìm số thực $x \ge 0$, biết $\sqrt{x} = 7$: \shortans{49}\par\vspace{0.8mm}
\noindent\textbf{Câu 9 (0,25 điểm).} Nghiệm nguyên dương của phương trình $(x - 4)(x + 3) = 0$ là: \shortans{4}\par\vspace{0.8mm}
\noindent\textbf{Câu 10 (0,25 điểm).} Nghiệm của phương trình $\sqrt{2x + 5} = 5$ là $x =$ \shortans{10}\par\vspace{0.8mm}
\vspace{0.8mm}
\noindent\textbf{Câu 11 (2,5 điểm). Giải các hệ phương trình sau:}
\vspace{0.6mm}

\noindent
\begin{tabularx}{\linewidth}{|X|X|}
\hline
\oBaiTap[21]{a) (1,25 điểm)}{$\begin{cases} 2x + y = 8 \\ x - y = 1 \end{cases}$} & 
\oBaiTap[21]{b) (1,25 điểm)}{$\begin{cases} 2(x - y) + 3y = 7 \\ 3x + 2y = 11 \end{cases}$} \\
\hline
\end{tabularx}

\newpage
% =============================================================================
% ========================= TOÁN 9 - TUẦN 06 - BUỔI 01: MÃ ĐỀ 103 =============
% ========================= MẶT SAU (TRANG 2/2) ================================
% =============================================================================

\noindent
\begin{tabularx}{\linewidth}{|X|c|}
\hline
\textbf{TOÁN 9 - TUẦN 06 - BUỔI 01} & \textbf{Mã đề: 103 - Trang 2/2} \\
\hline
\end{tabularx}

\vspace{1.2mm}

\noindent\textbf{Câu 12 (2,5 điểm). Biến đổi và rút gọn biểu thức:}
\vspace{0.6mm}

\noindent
\begin{tabularx}{\linewidth}{|X|X|}
\hline
\oBaiTap[13]{a) (1,25 điểm)}{Rút gọn với $x > 0, x \ne 16$:\par $A = \dfrac{\sqrt{x}}{\sqrt{x} - 4} - \dfrac{16}{x - 4\sqrt{x}}$} & 
\oBaiTap[13]{b) (1,25 điểm)}{Rút gọn với $x > 0, x \ne 1$:\par $B = \left(\dfrac{1}{\sqrt{x} - 1} + \dfrac{1}{\sqrt{x} + 1}\right) \cdot \dfrac{x - 1}{2\sqrt{x}}$} \\
\hline
\end{tabularx}

\vspace{2.5mm}

\noindent
\begin{tabularx}{\linewidth}{|X|}
\hline
\vspace{0.5mm}
\textbf{Câu 13 (2,5 điểm).} Hai bạn An và Bình đến một nhà sách mua đồ dùng học tập. Bạn An mua $2$ quyển vở và $3$ chiếc bút bi cùng loại hết tổng cộng $31\,000$ đồng. Bạn Bình mua $5$ quyển vở và $2$ chiếc bút bi cùng loại hết tổng cộng $50\,000$ đồng. Tính giá tiền của mỗi quyển vở và mỗi chiếc bút bi.\par
\vspace{0.4mm}
\noindent\textit{Bài giải:}\dongke{22}
\\ \hline
\end{tabularx}

\newpage
% =============================================================================
% ========================= TOÁN 9 - TUẦN 06 - BUỔI 01: MÃ ĐỀ 104 =============
% ========================= MẶT TRƯỚC (TRANG 1/2) ==============================
% =============================================================================

\noindent
\begin{tabularx}{\linewidth}{|X|c|}
\hline
\textbf{PHIẾU KIỂM TRA ĐỊNH KỲ TOÁN 9 - TUẦN 06 - BUỔI 01} & \textbf{Mã đề: 104 - Trang 1/2} \\
\hline
\end{tabularx}

\vspace{1.0mm}

\noindent
\begin{tabularx}{\linewidth}{|p{3.2cm}|X|}
\hline
\centering\textbf{ĐIỂM SỐ} & \centering\textbf{LỜI NHẬN XÉT CỦA THẦY CÔ} \tabularnewline
\hline
\rule{0pt}{1.8cm} & \tabularnewline
\hline
\end{tabularx}

\vspace{1.2mm}
\noindent\textbf{Câu 1 (0,25 điểm).} Cho $\Delta ABC$ vuông tại $A$ có $AB = 8, BC = 10$. Khi đó $\sin B =$ \shortans{\dfrac{3}{5}}\par\vspace{0.8mm}
\noindent\textbf{Câu 2 (0,25 điểm).} Cho $\Delta ABC$ vuông tại $A$ có $AB = 8, BC = 10$. Khi đó $\cos B =$ \shortans{\dfrac{4}{5}}\par\vspace{0.8mm}
\noindent\textbf{Câu 3 (0,25 điểm).} Cho $\Delta ABC$ vuông tại $A$ có $AB = 3, AC = 4$. Khi đó $\tan B =$ \shortans{\dfrac{4}{3}}\par\vspace{0.8mm}
\noindent\textbf{Câu 4 (0,25 điểm).} Cho $\Delta ABC$ vuông tại $A$, đường cao $AH$. Biết $BH=3, CH=12$. Độ dài $AH$ là: \shortans{6}\par\vspace{0.8mm}
\noindent\textbf{Câu 5 (0,25 điểm).} Một ngọn tháp cao $10\text{ m}$ có bóng trên mặt đất dài $10\text{ m}$. Góc tạo bởi tia nắng và mặt đất là: \shortans{45^\circ}\par\vspace{0.8mm}
\noindent\textbf{Câu 6 (0,25 điểm).} Cặp $(2; y)$ là nghiệm của phương trình $2x + y = 6$. Khi đó $y =$ \shortans{2}\par\vspace{0.8mm}
\noindent\textbf{Câu 7 (0,25 điểm).} Tìm $m$ để cặp $(2; 2)$ là nghiệm phương trình $mx + y = 6$: \shortans{2}\par\vspace{0.8mm}
\noindent\textbf{Câu 8 (0,25 điểm).} Tìm số thực $x \ge 0$, biết $\sqrt{x} = 4$: \shortans{16}\par\vspace{0.8mm}
\noindent\textbf{Câu 9 (0,25 điểm).} Nghiệm nguyên dương của phương trình $(x - 6)(x + 2) = 0$ là: \shortans{6}\par\vspace{0.8mm}
\noindent\textbf{Câu 10 (0,25 điểm).} Nghiệm của phương trình $\sqrt{4x + 9} = 5$ là $x =$ \shortans{4}\par\vspace{0.8mm}
\vspace{0.8mm}
\noindent\textbf{Câu 11 (2,5 điểm). Giải các hệ phương trình sau:}
\vspace{0.6mm}

\noindent
\begin{tabularx}{\linewidth}{|X|X|}
\hline
\oBaiTap[21]{a) (1,25 điểm)}{$\begin{cases} 3x + y = 10 \\ x - y = 2 \end{cases}$} & 
\oBaiTap[21]{b) (1,25 điểm)}{$\begin{cases} 3(x - 1) + 2y = 7 \\ 2x + y = 6 \end{cases}$} \\
\hline
\end{tabularx}

\newpage
% =============================================================================
% ========================= TOÁN 9 - TUẦN 06 - BUỔI 01: MÃ ĐỀ 104 =============
% ========================= MẶT SAU (TRANG 2/2) ================================
% =============================================================================

\noindent
\begin{tabularx}{\linewidth}{|X|c|}
\hline
\textbf{TOÁN 9 - TUẦN 06 - BUỔI 01} & \textbf{Mã đề: 104 - Trang 2/2} \\
\hline
\end{tabularx}

\vspace{1.2mm}

\noindent\textbf{Câu 12 (2,5 điểm). Biến đổi và rút gọn biểu thức:}
\vspace{0.6mm}

\noindent
\begin{tabularx}{\linewidth}{|X|X|}
\hline
\oBaiTap[13]{a) (1,25 điểm)}{Rút gọn với $x > 0, x \ne 25$:\par $A = \dfrac{\sqrt{x}}{\sqrt{x} - 5} - \dfrac{25}{x - 5\sqrt{x}}$} & 
\oBaiTap[13]{b) (1,25 điểm)}{Rút gọn với $x > 0, x \ne 16$:\par $B = \left(\dfrac{1}{\sqrt{x} - 4} + \dfrac{1}{\sqrt{x} + 4}\right) \cdot \dfrac{x - 16}{2\sqrt{x}}$} \\
\hline
\end{tabularx}

\vspace{2.5mm}

\noindent
\begin{tabularx}{\linewidth}{|X|}
\hline
\vspace{0.5mm}
\textbf{Câu 13 (2,5 điểm).} Hai bạn An và Bình đến một nhà sách mua đồ dùng học tập. Bạn An mua $3$ quyển vở và $4$ chiếc bút bi cùng loại hết tổng cộng $44\,000$ đồng. Bạn Bình mua $4$ quyển vở và $2$ chiếc bút bi cùng loại hết tổng cộng $42\,000$ đồng. Tính giá tiền của mỗi quyển vở và mỗi chiếc bút bi.\par
\vspace{0.4mm}
\noindent\textit{Bài giải:}\dongke{22}
\\ \hline
\end{tabularx}

\end{document}
